\begin{figure}[p]
\centering
\begin{tikzpicture}[x=1cm, y=1cm, font=\small,
  node distance=0.6cm,
  block/.style={draw=engblue, fill=engblue!8, thick, minimum width=2.4cm,
                minimum height=1.0cm, align=center, rounded corners=2pt},
  noisesrc/.style={draw=engred, fill=engred!8, thick, minimum width=2.0cm,
                   minimum height=0.7cm, align=center, rounded corners=2pt,
                   font=\footnotesize},
  arrow/.style={->, thick, engblue},
  noisearrow/.style={->, thick, engred, dashed}
]

%% Panel A: the measurement chain ===========================================
\node[anchor=west, font=\bfseries] at (-0.5, 9.6) {A. The measurement chain};

\node[block] (phenomenon) at (0.5, 8) {Physical\\phenomenon};
\node[block] (transducer) at (3.5, 8) {Transducer\\(thermistor, gauge,\\photodiode, ...)};
\node[block] (cond)       at (6.7, 8) {Signal\\conditioning\\(gain, filter)};
\node[block] (adc)        at (9.8, 8) {ADC\\(quantise,\\sample)};
\node[block] (record)     at (12.9, 8) {Digital\\record};

\draw[arrow] (phenomenon) -- (transducer) node[midway, above, font=\scriptsize]{$T,\,F,\,P$};
\draw[arrow] (transducer) -- (cond)       node[midway, above, font=\scriptsize]{$R,\,V,\,I$};
\draw[arrow] (cond)       -- (adc)        node[midway, above, font=\scriptsize]{$V_{\text{cond}}$};
\draw[arrow] (adc)        -- (record)     node[midway, above, font=\scriptsize]{count};

% Noise sources injected at each interface
\node[noisesrc] (n1) at (3.5, 6.3) {self-heating,\\thermal contact};
\node[noisesrc] (n2) at (6.7, 6.3) {Johnson noise,\\drift, EMI};
\node[noisesrc] (n3) at (9.8, 6.3) {quantisation,\\$V_{\text{ref}}$ drift};
\node[noisesrc] (n4) at (12.9, 6.3) {clock jitter,\\packet loss};

\draw[noisearrow] (n1) -- (transducer);
\draw[noisearrow] (n2) -- (cond);
\draw[noisearrow] (n3) -- (adc);
\draw[noisearrow] (n4) -- (record);

\node[anchor=west, font=\scriptsize\itshape, color=engred] at (-0.5, 5.4)
  {Each interface adds an uncertainty term. The transport archetype is the accounting.};

%% Panel B: the calibration loop ===========================================
\node[anchor=west, font=\bfseries] at (-0.5, 4.4) {B. The calibration loop};

\node[block, minimum width=2.6cm] (sensor) at (1.4, 3) {Sensor\\under test};
\node[block, minimum width=2.6cm] (refstd) at (5.6, 3) {Reference\\standard};
\node[block, minimum width=2.6cm] (fit)    at (9.8, 3) {Fit\\calibration\\function};
\node[block, minimum width=2.6cm] (record2) at (13.8, 3)
  {Record\\$\hat{f}$, $u(\hat{f})$};

\draw[arrow] (sensor) -- (refstd)
  node[midway, above, font=\scriptsize]{paired readings};
\draw[arrow] (refstd) -- (fit)
  node[midway, above, font=\scriptsize]{$\{(x_i, y_i)\}$};
\draw[arrow] (fit) -- (record2)
  node[midway, above, font=\scriptsize]{coeffs};

% feedback loop
\draw[arrow, rounded corners=4pt]
  (record2.south) -- ++(0, -0.7) -| (sensor.south);
\node[anchor=north, font=\scriptsize, color=engblue] at (7.4, 1.7)
  {periodic recalibration};

%% Panel C: spec sheet vs working uncertainty ==============================
\node[anchor=west, font=\bfseries] at (-0.5, 0.7)
  {C. Datasheet specification vs working uncertainty};

% Bars: stacked uncertainty contributions
\foreach \y/\label in {-0.5/{datasheet}, -1.5/{working}} {
  \node[anchor=east, font=\footnotesize] at (1.5, \y) {\label};
}

% datasheet: one segment
\fill[engblue!30] (1.7, -0.85) rectangle (4.0, -0.15);
\draw (1.7, -0.85) rectangle (4.0, -0.15);
\node[font=\scriptsize] at (2.85, -0.5) {manufacturer $\pm$};

% working: many segments
\definecolor{seg1}{RGB}{31, 78, 121}
\definecolor{seg2}{RGB}{86, 130, 168}
\definecolor{seg3}{RGB}{142, 178, 207}
\definecolor{seg4}{RGB}{200, 100, 80}
\definecolor{seg5}{RGB}{153, 42, 42}
\definecolor{seg6}{RGB}{120, 120, 120}

\fill[seg1!50] (1.7, -1.85) rectangle (3.0, -1.15);
\fill[seg2!60] (3.0, -1.85) rectangle (4.2, -1.15);
\fill[seg3!70] (4.2, -1.85) rectangle (5.4, -1.15);
\fill[seg4!50] (5.4, -1.85) rectangle (6.6, -1.15);
\fill[seg5!60] (6.6, -1.85) rectangle (8.0, -1.15);
\fill[seg6!50] (8.0, -1.85) rectangle (9.6, -1.15);
\draw (1.7, -1.85) rectangle (9.6, -1.15);

\node[font=\tiny] at (2.35, -1.5) {sensor};
\node[font=\tiny] at (3.6,  -1.5) {ref.\ res.};
\node[font=\tiny] at (4.8,  -1.5) {$V_{\text{ref}}$};
\node[font=\tiny] at (6.0,  -1.5) {quant.};
\node[font=\tiny] at (7.3,  -1.5) {self-heat};
\node[font=\tiny] at (8.8,  -1.5) {drift};

\node[anchor=west, font=\scriptsize\itshape, color=engred] at (10.0, -1.5)
  {The datasheet $\pm$ is one term in a longer budget.};

% axis
\draw[->, thick] (1.7, -2.3) -- (10.2, -2.3) node[right, font=\scriptsize]
  {combined standard uncertainty $u_c$};

\end{tikzpicture}
\caption[The sensor measurement chain, the calibration loop, and the
working uncertainty budget]{Three views of a sensor in working use.
Panel A: the measurement chain from physical phenomenon to digital
record, with the dominant noise or drift source at each interface.
Panel B: the calibration loop closes through a reference standard,
fits a calibration function, records the function and its
uncertainty, and is reopened at the next recalibration interval.
Panel C: the manufacturer's specification on the datasheet is one
term in the working uncertainty budget; the budget collects the
sensor's own residuals, the conditioning chain, the reference
resistor, the ADC reference, quantisation, self-heating, and
drift. A reader who reports the datasheet $\pm$ as if it were the
working uncertainty has substituted one term for the sum.}
\label{fig:vol01:ch05:sensor-chain}
\end{figure}
